\section{Experiments}

\subsection{Datasets}

To comprehensively evaluate the proposed Unified Dual-Mask Physical Model across diverse reflectance properties and scene complexities, we construct a multi-domain light field (LF) benchmark. The selected datasets encompass a wide spectrum of surface characteristics, ranging from ideal Lambertian to highly challenging non-Lambertian (e.g., specular, scattering, and transparent) materials. Specifically, our experiments are conducted on three primary datasets: the HCI New dataset [1], the UrbanLF-Syn dataset [2], and a curated Non-Lambertian dataset.

\paragraph{Dataset Descriptions.}

\textbf{HCI New Dataset.} The HCI New dataset [1] is a widely recognized synthetic benchmark primarily consisting of Lambertian scenes. It includes multiple complex scenes (e.g., \textit{boxes}, \textit{cotton}, \textit{dino}, and \textit{stratified}) with rich textures and intricate geometric structures. The dataset provides an angular resolution of $9 \times 9$ (81 views) and a spatial resolution of $512 \times 512$ pixels. The ground truth depth maps are perfectly generated by the rendering engine, providing pixel-accurate supervision. We utilize the official training and validation splits to assess the model's foundational depth estimation capability under ideal Lambertian assumptions.

\textbf{UrbanLF-Syn Dataset.} To evaluate the model's performance in mixed and complex urban environments, we employ the UrbanLF-Syn dataset [2]. This dataset comprises 170 synthetic scenes featuring diverse urban structures, varying illumination conditions, and mixed surface reflectance (i.e., a combination of Lambertian and mild non-Lambertian regions). The angular and spatial resolutions are standardized to $9 \times 9$ and $512 \times 512$, respectively. The large number of scenes provides robust training data for the mixed/urban domain, and the depth ground truth is derived from the synthetic rendering pipeline.

\textbf{Non-Lambertian Dataset.} Evaluating depth estimation on non-Lambertian surfaces remains a formidable challenge due to the violation of the photo-consistency assumption in Epipolar Plane Images (EPIs). For this domain, we utilize a specialized Non-Lambertian dataset featuring severe specular highlights, scattering, and transparent/mirror (ToM) surfaces. The scenes are synthesized utilizing measured Bidirectional Reflectance Distribution Functions (BRDFs) from the MERL BRDF database [3] to ensure physically accurate complex reflectance. Crucially, this dataset suffers from severe data scarcity, containing only 4 training scenes. The angular and spatial resolutions are $9 \times 9$ and $512 \times 512$. The extreme scarcity of non-Lambertian LF data constitutes a significant domain imbalance, serving as a rigorous testbed for the generalization and physical modeling capabilities of our unified framework.

\paragraph{Dataset Selection Rationale and Challenges.}

The selection of these three datasets is deliberately designed to cover the full trajectory of light field depth estimation challenges. The \textit{HCI New} dataset serves as the baseline for evaluating the upper bound of performance under standard Lambertian conditions. The \textit{UrbanLF-Syn} dataset introduces real-world complexity and mixed materials, testing the model's robustness to moderate domain shifts and occlusions. 

The most critical challenge lies in the \textit{Non-Lambertian} dataset. In non-Lambertian regions, the fundamental EPI linearity assumption breaks down, rendering traditional and even many deep learning-based EPI architectures ineffective. Furthermore, the extreme data scarcity (only 4 scenes) creates a severe long-tail distribution across domains. This intentional inclusion allows us to rigorously validate whether the proposed Dual-Mask Physical Model and Angular FFT Material Classification can effectively decouple domain-specific physical features and prevent the model from overfitting to the dominant Lambertian/Mixed domains.

\paragraph{Data Preprocessing and Sampling Strategy.}

To ensure consistency across the multi-domain benchmark, all light field images are preprocessed to a uniform spatial resolution of $512 \times 512$ and an angular resolution of $9 \times 9$. Depth maps are normalized to a unified range to facilitate stable joint optimization of the depth regression and domain classification objectives.

To mitigate the severe domain imbalance—particularly the drastic disparity between the 170 UrbanLF scenes and the 4 Non-Lambertian scenes—we implement a \textbf{Domain-Balanced Sampling} strategy during training. Specifically, we employ a `WeightedRandomSampler` that dynamically assigns higher sampling probabilities to the minority domains (Non-Lambertian) and lower probabilities to the majority domains (UrbanLF-Syn). This ensures that each training batch contains a balanced representation of all physical domains, preventing the network from being biased toward the Lambertian/Mixed priors and encouraging the dual-mask modules to learn domain-invariant and domain-specific physical features effectively. The datasets are strictly partitioned into training and validation sets, with no overlap in scene identities to ensure a fair evaluation of generalization.

\paragraph{Table I: Summary of Datasets.}

\begin{table*}[!t]
\small
\caption{Dataset Scene Type No. of Scenes.}
\label{tab:table1}
\centering
\resizebox{\textwidth}{!}{
\begin{tabular*}{\textwidth}{@{\extracolsep{\fill}}lllllll}
\toprule
Dataset \& Scene Type \& No. of Scenes \& Angular Res. \& Spatial Res. \& Depth Source \& Key Characteristics \\
\midrule
\textbf{HCI New} \& Lambertian \& 4 train + 4 val \& $9 \times 9$ \& $512 \times 512$ \& Synthetic GT \& Ideal photo-consistency, rich textures, standard benchmark. \\
\textbf{UrbanLF-Syn} \& Mixed / Urban \& 170 \& $9 \times 9$ \& $512 \times 512$ \& Synthetic GT \& Complex urban structures, mixed reflectance, large scale. \\
\textbf{Non-Lambertian} \& Non-Lambertian \& 4 \& $9 \times 9$ \& $512 \times 512$ \& Synthetic GT (MERL BRDF) \& Severe specular/scattering, EPI assumption breakdown, extreme data scarcity. \\
\bottomrule
\end{tabular*}
}
\end{table*}
\textit{\}Note: The "No. of Scenes" for HCI New refers to the specific subsets (e.g., training and stratified) used in our experiments. All datasets are strictly split into training and validation sets with no overlap.*

\subsubsection{Implementation Details}

\paragraph{Hardware and Software Environment.}
All experiments were implemented using the PyTorch 2.0 framework and trained on a high-performance computing cluster equipped with four NVIDIA RTX A6000 GPUs (48GB VRAM each). The software environment was built on CUDA 11.8 and cuDNN 8.6 to accelerate tensor operations. To maximize hardware utilization and synchronize gradients across multiple devices, Distributed Data Parallel (DDP) training was employed. 

\paragraph{Training Hyperparameters and Optimization.}
The proposed Unified Dual-Mask Physical Model is built upon the EPINet4Dir V3 backbone, which extracts 4-directional Epipolar Plane Images (EPIs). We employ the AdamW optimizer with a weight decay of $1 \times 10^{-4}$ to mitigate overfitting, particularly in the data-scarce Non-Lambertian domain. The initial learning rate is set to $1 \times 10^{-4}$, which is dynamically adjusted using a Cosine Annealing learning rate scheduler with a minimum learning rate of $1 \times 10^{-6}$ and a warm-up phase of 5 epochs. 

Due to the substantial memory footprint of 4D light field data (e.g., $9 \times 9$ angular views at $512 \times 512$ spatial resolution), the per-GPU batch size is restricted to 4. To maintain a sufficient effective batch size for stable gradient estimation, we utilize gradient accumulation with an accumulation step of 4, yielding an effective batch size of 64 across the four GPUs. The model is trained for a maximum of 300 epochs. Early stopping with a patience of 30 epochs is applied based on the validation Mean Absolute Error (MAE) to prevent overfitting. The key training hyperparameters are summarized in Table II.

\textbf{TABLE II}
\textbf{SUMMARY OF KEY TRAINING HYPERPARAMETERS}

\begin{table}[!t]
\caption{Hyperparameter Configuration Hyperparameter.}
\label{tab:table2}
\centering
\begin{tabular*}{\linewidth}{@{\extracolsep{\fill}}llll}
\toprule
Hyperparameter \& Configuration \& Hyperparameter \& Configuration \\
\midrule
\textbf{Optimizer} \& AdamW \& \textbf{Base Model} \& EPINet4Dir V3 \\
\textbf{Initial Learning Rate} \& $1 \times 10^{-4}$ \& \textbf{Weight Decay} \& $1 \times 10^{-4}$ \\
\textbf{Min Learning Rate} \& $1 \times 10^{-6}$ \& \textbf{Warm-up Epochs} \& 5 \\
\textbf{Batch Size (per GPU)} \& 4 \& \textbf{Gradient Accum.} \& 4 \\
\textbf{Effective Batch Size} \& 64 \& \textbf{Max Epochs} \& 300 \\
\textbf{Early Stopping} \& Patience 30 \& \textbf{Scheduler} \& Cosine Annealing \\
\bottomrule
\end{tabular*}
\end{table}
Building upon these comprehensive empirical validations, we now distill our core contributions and outline the broader impact of this work.